% Generated draft: requires mathematical and source review.
\documentclass[11pt]{article}
\usepackage[margin=25mm]{geometry}
\usepackage{fontspec}
\setmainfont{Latin Modern Roman}
\usepackage{amsmath,amssymb,mathtools,unicode-math}
\setmathfont{Latin Modern Math}
\usepackage{enumitem,needspace,xurl,hyperref}
\setlength{\parindent}{0pt}
\setlength{\parskip}{5pt}
\setlength{\emergencystretch}{3em}
\newcommand{\N}{\mathbb{N}}
\newcommand{\Z}{\mathbb{Z}}
\newcommand{\Q}{\mathbb{Q}}
\newcommand{\R}{\mathbb{R}}
\newcommand{\C}{\mathbb{C}}
\newcommand{\F}{\mathbb{F}}
\newcommand{\A}{\mathbb{A}}
\newcommand{\PP}{\mathbb{P}}
\newcommand{\E}{\mathbb{E}}
\newcommand{\eps}{\varepsilon}
\newcommand{\dd}{\,\mathrm{d}}
\begin{document}
\section*{UnsolvedMath ID: 10900021}
\textbf{Source problem code:} catalogue entry
\subsection{Definitions and mathematical statement}
A finitely generated group \(G\) is hyperbolic if a Cayley graph has uniformly thin geodesic triangles, and freely indecomposable if it has no decomposition \(G\cong A*B\) with both factors nontrivial. Under the standard convention, cubulated means admitting a proper cocompact action on a CAT\((0)\) cube complex.

Futer asks whether every freely indecomposable cubulated hyperbolic group contains a closed hyperbolic surface subgroup. Using orientable surfaces, the requested conclusion is
\[
 \exists g\geq2\quad\exists\text{ injective homomorphism}
 \quad\pi_1(\Sigma_g)\hookrightarrow G,
\]
where \(\Sigma_g\) is the closed connected orientable surface of genus \(g\). No quasiconvexity, finite-index or geometric embedding condition is imposed on the subgroup.

As printed, the hypotheses do not exclude elementary groups: for example, the infinite cyclic group is freely indecomposable, hyperbolic and cubulated under these definitions, but cannot contain a nonabelian closed hyperbolic surface group. Thus a non-elementary/one-ended or other intended exclusion is needed for the usual conjectural target; it is not silently added here. The background identifies this as the cubulated case of Gromov's Surface Subgroup Conjecture, but also supplies no explicit exclusion. The intended convention for cubulation and the treatment of torsion must be recovered from the source context.
\subsection{Short English statement}
Ask whether the intended nontrivial class of freely indecomposable cubulated hyperbolic groups always contains a closed hyperbolic surface group. The printed hypotheses omit elementary-group exclusions.
\subsection{Sources}
\begin{itemize}
\item Futer, attributed item 4.2 in Delp, Hoffoss and Manning, \emph{Problems In Groups, Geometry, and Three-Manifolds} (2015), supplied PDF page 3. \url{https://arxiv.org/pdf/1512.04620#page=3}; status link \url{https://arxiv.org/abs/1512.04620}. Recorded access 2026-07-29; extraction source-tex checked against rendered PDF text, as reported in the record.
\item The supplied summary attributes the hyperbolic \(3\)-manifold case to Kahn--Markovic and lists ``Immersed essential surfaces in hyperbolic 3-manifolds,'' Duke 2012, mentioning good pants. These are the supplied bibliographic details, not independently verified; the precise manifold hypotheses and citation require checking.
\item Local UnsolvedMath registry, \texttt{data/unsolvedmath/problems.json}, permanent ID \(10900021\), exact source code \texttt{catalogue entry}; supplied OPEN-TRIAGE status and general-case open report.
\end{itemize}
\end{document}
